\chapter{Complexity — MARM Compile}
%%% generated by the blueprint pipeline from TCSlib.Complexity.SpaceComplexity.Machines.MARMCompile
\section{Overview}

Declarations of the Lean module \texttt{TCSlib.Complexity.SpaceComplexity.Machines.MARMCompile}.

\section{Declarations}

\begin{definition}[Memory tape as a partial tape]
\lean{Complexity.LogProg.memTape}
\label{Complexity.LogProg.memTape}
\leanok
Given a Boolean-valued memory function $M$ on the integer-indexed cells, represent it as
a partial tape that contains the symbol $1$ at cells where $M$ is true and is blank at
cells where $M$ is false.
\end{definition}

\begin{definition}[Adding a blank work tape]
\lean{Complexity.LogProg.liftAct}
\label{Complexity.LogProg.liftAct}
\leanok
\uses{Turing.Action}
Given an action of a Turing machine with $K$ work tapes, form an action with $K+1$ work
tapes by retaining its input-head movement, output, and successor state, preserving the
actions on its first $K$ work tapes, and adding a final work tape that writes nothing
and does not move its head.
\end{definition}

\begin{definition}[Configuration with an added memory tape]
\lean{Complexity.LogProg.memCfg}
\label{Complexity.LogProg.memCfg}
\leanok
\uses{Turing.Cfg}
Let $c$ be a configuration of a machine with $K$ work tapes on input $x$, let
$M\colon\mathbb Z\to\mathrm{Option}(\{0,1\})$ describe an additional tape, and let
$p\in\mathbb Z$. The resulting configuration has $K+1$ work tapes: it preserves the
state, input-head position, original work tapes and their head positions, and output of
$c$, and uses $M$ and $p$ as the contents and head position of the added tape.
\end{definition}

\begin{lemma}[Existing tape symbols after adjoining a memory tape]
\lean{Complexity.LogProg.memCfg_symbols_castSucc}
\label{Complexity.LogProg.memCfg_symbols_castSucc}
\leanok
\uses{Complexity.LogProg.memCfg, Turing.Cfg, Turing.Cfg.workTapeSymbols}
\difficulty{1}
Let $c$ be a configuration of a $K$-work-tape Turing machine, and adjoin a further tape
with contents $M$ and head position $p$. For every original tape index $i<K$, the symbol
under its head in the enlarged configuration is the same as the symbol under its head in
$c$.
\end{lemma}

\begin{lemma}[Symbol under the added tape head]
\lean{Complexity.LogProg.memCfg_symbols_last}
\label{Complexity.LogProg.memCfg_symbols_last}
\leanok
\uses{Complexity.LogProg.memCfg, Turing.Cfg, Turing.Cfg.workTapeSymbols}
\difficulty{1}
Let $c$ be a configuration of a $K$-work-tape Turing machine, and adjoin a tape with
contents $M$ and head position $p$. The symbol under the head of this added tape is
$M(p)$.
\end{lemma}

\begin{lemma}[Lifting a machine action]
\lean{Complexity.LogProg.liftAct_apply}
\label{Complexity.LogProg.liftAct_apply}
\leanok
\uses{Complexity.LogProg.liftAct, Complexity.LogProg.memCfg, Turing.Action, Turing.Action.apply, Turing.Cfg}
\difficulty{4}
Let an action of a $K$-work-tape machine act on a configuration with fixed input, and
extend the configuration by one work tape whose contents are a function $M\colon\mathbb
Z\to\mathrm{Option}(\mathrm{Bool})$ and whose head is at position $p$. Extend the action
by leaving this additional tape unwritten and stationary. Applying the extended action
to the extended configuration gives the extension, by the same $M$ and $p$, of the
configuration obtained by applying the original action.
\end{lemma}

\begin{definition}[Memory instruction selection]
\lean{Complexity.LogProg.mop}
\label{Complexity.LogProg.mop}
\leanok
\uses{Complexity.LogProg.CSt, Complexity.LogProg.MARM, Complexity.LogProg.MIns, Complexity.LogProg.MIns.isMem, Complexity.LogProg.Ph}
Let $A$ assign an instruction to each program node. For a state consisting of a program
node and a phase, the selector returns the instruction assigned at that node exactly
when the phase is the start phase and the assigned instruction is a memory instruction;
in every other state it returns no instruction.
\end{definition}

\begin{lemma}[Characterization of selected memory instructions]
\lean{Complexity.LogProg.mop_spec}
\label{Complexity.LogProg.mop_spec}
\leanok
\uses{Complexity.LogProg.CSt, Complexity.LogProg.MARM, Complexity.LogProg.MIns, Complexity.LogProg.MIns.isMem, Complexity.LogProg.Ph, Complexity.LogProg.mop}
\difficulty{4}
Let $A$ be a program with memory instructions, and let $q$ be a compiled-machine state.
If the memory-operation selector returns instruction $i$ at $q$, then there is a program
state $l$ such that $q$ is the program-running state at $l$ with its phase at the start,
$A$ assigns instruction $i$ to $l$, and $i$ is a memory instruction.
\end{lemma}

\begin{definition}[Actions for memory instructions]
\lean{Complexity.LogProg.memAct}
\label{Complexity.LogProg.memAct}
\leanok
\uses{Complexity.LogProg.CSt, Complexity.LogProg.MIns, Complexity.LogProg.Ph, Turing.Action}
Let $i$ be a memory instruction and let $r$ be the optional bit read from the cell under
the memory head. A register instruction gives an action with no head movement, tape
writes, output, or successor state. A read instruction gives no movement, writes, or
output and continues at the first successor if $r$ is true, and at the second otherwise.
A set instruction writes $1$ on the added tape and continues at its specified node; a
left or right instruction moves that tape's head one cell in the indicated direction and
continues at its specified node. In these last three cases, the input head and all other
work tapes stay put, and no output is emitted.
\end{definition}

\begin{definition}[Compilation with an added memory tape]
\lean{Complexity.LogProg.mcompileTM}
\label{Complexity.LogProg.mcompileTM}
\leanok
\uses{Complexity.LogProg.CSt, Complexity.LogProg.MARM, Complexity.LogProg.MARM.detPart, Complexity.LogProg.Ph, Complexity.LogProg.armProg, Complexity.LogProg.ctr, Complexity.LogProg.liftAct, Complexity.LogProg.memAct, Complexity.LogProg.mop, Turing.MultiTapeTM}
Let $A$ assign memory-capable instructions to program nodes, let $l_0$ be an initial
node, let $D$ be a Boolean-alphabet machine with $k_D$ work tapes and state set $SD$,
and let $q_0$ assign an initial state of $D$ to each of $d$ decider indices. The
compiled machine has $m+k_D+1$ work tapes and states that record a program node and
phase together with the decider's simulation data. It starts at $(l_0,\mathrm{start})$;
at each step it executes a memory instruction selected by the memory selector when one
is available, and otherwise uses the compiled register-program and decider transition.
\end{definition}

\begin{definition}[Finite machine produced by compilation]
\lean{Complexity.LogProg.mcompileFinTM}
\label{Complexity.LogProg.mcompileFinTM}
\leanok
\uses{Complexity.LogProg.CSt, Complexity.LogProg.MARM, Complexity.LogProg.Ph, Complexity.LogProg.mcompileTM, Turing.FinTM, Turing.MultiTapeTM}
Let $m$, $d$, and $k_D$ be natural numbers, let $Lambda$ and $SD$ be finite types with
decidable equality, and let $A$ assign to each element of $Lambda$ an instruction with
$m$ registers, $d$ decider indices, and labels in $Lambda$. Choose an initial label
$l_0$, a Boolean-alphabet Turing machine $D$ with $k_D$ work tapes and state set $SD$,
and a map $q_0$ assigning a state of $D$ to each of the $d$ decider indices. The
compilation of these data is a finite Boolean-alphabet Turing machine with $m+k_D+1$
work tapes.
\end{definition}

\begin{definition}[Deterministic register-machine compilation]
\lean{Complexity.LogProg.detM}
\label{Complexity.LogProg.detM}
\leanok
\uses{Complexity.LogProg.CSt, Complexity.LogProg.MARM, Complexity.LogProg.MARM.detPart, Complexity.LogProg.Ph, Complexity.LogProg.armProg, Complexity.LogProg.compileTM, Turing.MultiTapeTM}
Let $A$ assign memory-capable instructions to program nodes, let $l_0$ be an initial
node, let $D$ be a Boolean-alphabet machine with $k_D$ work tapes and states $SD$, let
$q_0$ assign a state of $D$ to each of $d$ decider indices, and choose a Boolean
parameter $\beta$. The resulting machine is the compilation of the register program
obtained from $A$ using parameter $\beta$, together with $D$ and the given initial-state
assignment; it has $m+k_D$ work tapes and the compiled state set.
\end{definition}

\begin{lemma}[Call behavior is independent of memory-choice replacement]
\lean{Complexity.LogProg.armProg_mdetPart_call}
\label{Complexity.LogProg.armProg_mdetPart_call}
\leanok
\uses{Complexity.LogProg.MARM, Complexity.LogProg.MARM.detPart, Complexity.LogProg.MARM.detPart_det, Complexity.LogProg.MARM.detPart_mem, Complexity.LogProg.MIns.isMem, Complexity.LogProg.armProg, Complexity.LogProg.insCall}
\difficulty{4}
Let $A$ assign to each label either a deterministic instruction or a memory instruction,
and let $l_0\in\Lambda$ and $\beta,\beta'\in\{\mathrm{false},\mathrm{true}\}$. The
compiled register programs obtained by retaining deterministic instructions and
replacing every memory instruction with the corresponding return of $\beta$ and of
$\beta'$ have identical call maps.
\end{lemma}

\begin{lemma}[Deterministic transition independence]
\lean{Complexity.LogProg.detM_tr}
\label{Complexity.LogProg.detM_tr}
\leanok
\uses{Complexity.LogProg.CSt, Complexity.LogProg.MARM, Complexity.LogProg.MARM.detPart, Complexity.LogProg.MARM.detPart_det, Complexity.LogProg.MARM.detPart_mem, Complexity.LogProg.MIns.isMem, Complexity.LogProg.Ph, Complexity.LogProg.armProg_mdetPart_call, Complexity.LogProg.compileTM, Complexity.LogProg.ctr_congr, Complexity.LogProg.detM, Complexity.LogProg.insTr, Complexity.LogProg.mop, Turing.MultiTapeTM}
\difficulty{5}
Let $A$ be a memory-augmented register program, let $D$ be a decider, and choose an
initial decider state for each decider index. For any Boolean answer $eta$ and compiled
state $q$ at which the memory-choice map of $A$ is undefined, the transition of the
deterministic compilation with answer parameter $eta$ agrees with the transition using
answer parameter false.
\end{lemma}

\begin{lemma}[One-step simulation with an appended memory tape]
\lean{Complexity.LogProg.mcompileTM_step}
\label{Complexity.LogProg.mcompileTM_step}
\leanok
\uses{Complexity.LogProg.CSt, Complexity.LogProg.MARM, Complexity.LogProg.Ph, Complexity.LogProg.detM, Complexity.LogProg.detM_tr, Complexity.LogProg.liftAct_apply, Complexity.LogProg.mcompileTM, Complexity.LogProg.memCfg, Complexity.LogProg.mop, Turing.Cfg, Turing.Cfg.inputSymbol, Turing.Cfg.workTapeSymbols, Turing.MultiTapeTM, Turing.MultiTapeTM.step}
\difficulty{4}
Let $A$ be a memory-augmented register program, let $l_0$ be its start label, let $D$ be
a decider with initial-state map $q_0$, and let $\beta$ be a Boolean value. For a
configuration $c$ of the corresponding deterministic compiled machine, suppose that no
nonhalting state of $c$ enables a memory instruction of $A$. Then, for any memory-tape
contents $M$ and head position $p$, one step of the memory-augmented compiled machine
from the configuration obtained by appending that tape equals the configuration obtained
by taking one step of the deterministic compiled machine from $c$ and then appending the
unchanged memory tape and head position.
\end{lemma}

\begin{lemma}[Simulation with Unused Memory]
\lean{Complexity.LogProg.mrunFrom_eq}
\label{Complexity.LogProg.mrunFrom_eq}
\leanok
\uses{Complexity.LogProg.CSt, Complexity.LogProg.MARM, Complexity.LogProg.Ph, Complexity.LogProg.detM, Complexity.LogProg.mcompileTM, Complexity.LogProg.mcompileTM_step, Complexity.LogProg.memCfg, Complexity.LogProg.mop, Turing.Cfg, Turing.MultiTapeTM, Turing.MultiTapeTM.runFrom, Turing.MultiTapeTM.runFrom_succ_eq_step'}
\difficulty{4}
Let $A$ be a memory-augmented register program, let $l_0$ be its initial state, let $D$
be a decider machine, let $q_0$ assign an initial decider state to each decider, and fix
a Boolean value $\beta$. Let $c$ be a configuration of the deterministic compiled
machine, let $M:\mathbb Z\to\{0,1\}$ be a memory tape, let $p\in\mathbb Z$ be its head
position, and let $T$ be a natural number. Assume that at every time $t<T$, whenever the
deterministic run from $c$ is in a state $q$, the program has no memory instruction at
$q$. Then, for every $t\le T$, the compiled memory-augmented machine started from the
configuration obtained by adding memory tape $M$ with head at $p$ agrees with the
deterministic run at time $t$, with that same memory tape and head position attached.
\end{lemma}

\begin{lemma}[Halting step for memory operations]
\lean{Complexity.LogProg.detM_step_mem}
\label{Complexity.LogProg.detM_step_mem}
\leanok
\uses{Complexity.LogProg.CSt, Complexity.LogProg.MARM, Complexity.LogProg.MARM.detPart, Complexity.LogProg.MARM.detPart_mem, Complexity.LogProg.Ph, Complexity.LogProg.armProg, Complexity.LogProg.compileTM, Complexity.LogProg.ctr, Complexity.LogProg.detM, Complexity.LogProg.insCall, Complexity.LogProg.insTr, Complexity.LogProg.mop, Complexity.LogProg.mop_spec, Complexity.LogProg.retAct, Turing.Action.apply, Turing.Cfg, Turing.MultiTapeTM, Turing.MultiTapeTM.step}
\difficulty{5}
Let $A$ be a program of memory instructions, let $l_0$ be an initial program location,
let $D$ be a multi-tape Turing machine, let $q_0$ assign an initial state of $D$ to each
decider index, and let $eta$ be a Boolean value. For any input and any configuration
$c$ of the compiled machine, if $c$ has state $q$ and the memory-instruction map at $q$
is defined, then the machine halts after one step and its output is the output of $c$
with $eta$ appended.
\end{lemma}

\begin{lemma}[Choice-independent memory precondition]
\lean{Complexity.LogProg.pre_mdetPart}
\label{Complexity.LogProg.pre_mdetPart}
\leanok
\uses{Complexity.LogProg.AConf, Complexity.LogProg.MARM, Complexity.LogProg.MARM.detPart, Complexity.LogProg.MARM.detPart_det, Complexity.LogProg.MARM.detPart_mem, Complexity.LogProg.MIns.isMem, Complexity.LogProg.Pre, Turing.MultiTapeTM}
\difficulty{4}
For every memory program, oracle, decider, initial-state assignment, space bound, input,
two Boolean replacements for choice instructions, and configuration, the preconditions
of the two resulting deterministic programs agree.
\end{lemma}

\begin{lemma}[Deterministic instruction step is independent of halt answer]
\lean{Complexity.LogProg.astep_mdetPart}
\label{Complexity.LogProg.astep_mdetPart}
\leanok
\uses{Complexity.LogProg.Ins, Complexity.LogProg.MARM, Complexity.LogProg.MARM.detPart, Complexity.LogProg.MARM.detPart_det, Complexity.LogProg.astep}
\difficulty{1}
Let $m,d$ be natural numbers, let $\Lambda$ be a type, let $A$ be a memory arithmetic
register machine, and suppose its instruction at state $l$ is the deterministic
instruction $i$. For any oracle, input bit string, Boolean values $\beta,\beta'$,
register contents $v$, and result value $r$, the arithmetic-machine steps using the two
deterministic parts that return $\beta$ or $\beta'$ on non-deterministic instructions
agree when started at state $l$ with contents $v$ and result $r$.
\end{lemma}

\begin{definition}[Configuration with a memory tape at a seam]
\lean{Complexity.LogProg.mseam}
\label{Complexity.LogProg.mseam}
\leanok
\uses{Complexity.LogProg.CSt, Complexity.LogProg.Ph, Complexity.LogProg.aseam, Complexity.LogProg.memCfg, Complexity.LogProg.memTape, Complexity.LogProg.seam, Turing.Cfg}
Let $x$ be a Boolean input list, $l$ a program node, $v\colon\mathrm{Fin}(m)\to\mathbb
N$ a register assignment, $M\colon\mathbb Z\to\{0,1\}$ a memory tape, $p\in\mathbb Z$
its head position, and $out$ a Boolean output list. The construction forms a
configuration on input $x$ with $m+k_D+1$ work tapes by taking the seam configuration
associated with $x,l,v,out$ and adjoining a memory tape whose contents are the
optional-symbol encoding of $M$ and whose head is at $p$.
\end{definition}

\begin{lemma}[Initial configuration of the compiled machine]
\lean{Complexity.LogProg.mcompileTM_init}
\label{Complexity.LogProg.mcompileTM_init}
\leanok
\uses{Complexity.LogProg.MARM, Complexity.LogProg.aseam, Complexity.LogProg.mcompileTM, Complexity.LogProg.memCfg, Complexity.LogProg.memTape, Complexity.LogProg.mseam, Complexity.LogProg.seam, Turing.MultiTapeTM, Turing.MultiTapeTM.initCfg}
\difficulty{5}
For every memory-augmented register program $A$, start label $l_0$, decider $D$,
initial-state map $q_0$, and bit string $x$, the compiled machine's initial
configuration on input $x$ is the configuration with label $l_0$, all register values
zero, memory identically false, memory-head position zero, and empty output.
\end{lemma}

\begin{definition}[Bounds on work-tape head positions]
\lean{Complexity.LogProg.MBox}
\label{Complexity.LogProg.MBox}
\leanok
\uses{Turing.Cfg}
Let $W,B,P\in\mathbb N$ and let $C$ be a configuration with $m+k_D+1$ work tapes. The
first $m$ work-tape heads each lie between $-1$ and $W$, inclusive; each of the next
$k_D$ work-tape heads has absolute position at most $B$; and the final work-tape head
has absolute position at most $P$.
\end{definition}

\begin{lemma}[Bounds for the initial compiled configuration]
\lean{Complexity.LogProg.mbox_mseam}
\label{Complexity.LogProg.mbox_mseam}
\leanok
\uses{Complexity.LogProg.CSt, Complexity.LogProg.MBox, Complexity.LogProg.Ph, Complexity.LogProg.aseam, Complexity.LogProg.memCfg, Complexity.LogProg.mseam, Complexity.LogProg.seam}
\difficulty{3}
Let $W,B,P$ be natural numbers, and form the initial compiled configuration from a bit
string $x$, label $l$, zero-valued registers, Boolean memory contents $M$, memory-head
position $p$, and output word $out$. If $|p|\le P$, then this configuration has each
register head in $[-1,W]$, each decider-work-tape head of absolute position at most $B$,
and memory-head position of absolute value at most $P$.
\end{lemma}

\begin{lemma}[Bounds after appending a memory tape]
\lean{Complexity.LogProg.mbox_memCfg}
\label{Complexity.LogProg.mbox_memCfg}
\leanok
\uses{Complexity.LogProg.MBox, Complexity.LogProg.memCfg, Turing.Cfg}
\difficulty{3}
Let $c$ be a configuration of a machine with $m+k_D$ work tapes, and append a memory
tape with contents $M$ and head position $p$. If each of the first $m$ work-tape heads
of $c$ lies in $[-1,W]$, each of its next $k_D$ work-tape heads has absolute position at
most $B$, and $|p|\le P$, then the resulting configuration has those same respective
bounds on its first $m$ heads and next $k_D$ heads, and its appended memory head has
absolute position at most $P$.
\end{lemma}

\begin{theorem}[Live-step simulation of a deterministic instruction]
\lean{Complexity.LogProg.mdet_live}
\label{Complexity.LogProg.mdet_live}
\leanok
\uses{Complexity.LogProg.MARM, Complexity.LogProg.MARM.detPart, Complexity.LogProg.MBox, Complexity.LogProg.Pre, Complexity.LogProg.armProg, Complexity.LogProg.arm_step, Complexity.LogProg.aseam, Complexity.LogProg.astep, Complexity.LogProg.avoid_of_live, Complexity.LogProg.compile_correct_from, Complexity.LogProg.detM, Complexity.LogProg.detM_step_mem, Complexity.LogProg.mbox_memCfg, Complexity.LogProg.mcompileTM, Complexity.LogProg.memTape, Complexity.LogProg.mop, Complexity.LogProg.mrunFrom_eq, Complexity.LogProg.mseam, Complexity.LogProg.seam, Turing.MultiTapeTM, Turing.MultiTapeTM.runFrom}
\difficulty{5}
Let $m,d,k_D$ be nonnegative integers, let $\Lambda$ and $SD$ be state sets, and fix a
Boolean input $x$, an augmented register machine $A$ with $m$ registers and $d$
deciders, an initial label $l_0$, an oracle, a $k_D$-work-tape Turing machine $D$,
initial decider states $q_0$, bounds $B,W,P$, labels $l,l'$, register assignments
$v,v'$, a Boolean memory tape $M$ with head position $p$, and an output word $out$.
Suppose the ordinary execution from $(l,v)$ with no prior answer takes one step to
$(l',v')$ with no answer, the precondition holds at $(l,v)$, every register satisfies
$|\mathrm{bits}(v(r)+1)|\le W$, and $|p|\le P$. Then for some finite time the compiled
Turing machine, started with input $x$, registers $v$, memory tape $M$ at $p$, and
output $out$, reaches exactly the corresponding configuration with label $l'$, registers
$v'$, the same memory tape and head position, and the same output; throughout this run
every configuration satisfies the bounds $W,B,P$ encoded by $\mathrm{MBox}$.
\end{theorem}

\begin{theorem}[Halting simulation of a deterministic instruction]
\lean{Complexity.LogProg.mdet_halt}
\label{Complexity.LogProg.mdet_halt}
\leanok
\uses{Complexity.LogProg.Ins, Complexity.LogProg.MARM, Complexity.LogProg.MARM.detPart, Complexity.LogProg.MBox, Complexity.LogProg.Pre, Complexity.LogProg.armProg, Complexity.LogProg.arm_run, Complexity.LogProg.arun, Complexity.LogProg.aseam, Complexity.LogProg.astep, Complexity.LogProg.astep_mdetPart, Complexity.LogProg.avoid_of_halt, Complexity.LogProg.compile_correct_from, Complexity.LogProg.detM, Complexity.LogProg.detM_step_mem, Complexity.LogProg.mbox_memCfg, Complexity.LogProg.mcompileTM, Complexity.LogProg.memTape, Complexity.LogProg.mop, Complexity.LogProg.mrunFrom_eq, Complexity.LogProg.mseam, Complexity.LogProg.pre_mdetPart, Complexity.LogProg.seam, Complexity.LogProg.seam_state_none, Turing.MultiTapeTM, Turing.MultiTapeTM.runFrom}
\difficulty{5}
Let $m,d,k_D$ be nonnegative integers, let $\Lambda$ and $SD$ be state sets, and fix a
Boolean input $x$, an augmented register machine $A$ with $m$ registers and $d$
deciders, an initial label $l_0$, an oracle, a $k_D$-work-tape Turing machine $D$,
initial decider states $q_0$, bounds $B,W,P$, a label $l$, register contents $v$, a
Boolean memory tape $M$ with head position $p$, an output word $out$, and a bit $b_0$.
Suppose the instruction at $l$ is the deterministic instruction $i$, and the ordinary
register-machine execution from $(l,v)$ with no prior answer halts with unchanged
registers $v$ and answer $b_0$. Suppose also that the stated precondition holds at this
configuration, every register satisfies $|\mathrm{bits}(v(r)+1)|\le W$, and $|p|\le P$.
Then the compiled Turing machine, started with input $x$, register contents $v$, memory
tape $M$ at position $p$, and output $out$, reaches a halting configuration whose output
is $out$ followed by $b_0$, and every configuration up to that time satisfies the
register, decider-work-tape, and memory-head bounds $W,B,P$ encoded by $\mathrm{MBox}$.
\end{theorem}

\begin{lemma}[One-step correspondence for memory instructions]
\lean{Complexity.LogProg.mmem_step}
\label{Complexity.LogProg.mmem_step}
\leanok
\uses{Complexity.LogProg.CSt, Complexity.LogProg.MARM, Complexity.LogProg.MIns.isMem, Complexity.LogProg.Ph, Complexity.LogProg.aseam, Complexity.LogProg.mcompileTM, Complexity.LogProg.memAct, Complexity.LogProg.memCfg, Complexity.LogProg.memTape, Complexity.LogProg.mop, Complexity.LogProg.mseam, Complexity.LogProg.mstep, Complexity.LogProg.seam, Turing.Action, Turing.Action.apply, Turing.Cfg.workTapeSymbols, Turing.MultiTapeTM, Turing.MultiTapeTM.step}
\difficulty{6}
Let $A$ be a program with memory instructions, and suppose its instruction at state $l$
is a memory instruction. For any oracle, input word, decider, initial decider states,
register values, Boolean memory tape, memory-head position, and output word, there are a
state $l'$, a Boolean memory tape $M'$, and a memory-head position $p'$ such that one
program step from $(l,v,\mathrm{none})$ with memory $(M,p)$ gives
$((l',v,\mathrm{none}),(M',p'))$, and one compiled-machine step from the corresponding
seam configuration gives the seam configuration for $l',v,M',p'$, with the output
unchanged.
\end{lemma}
